% Standalone problem document, generated from the adjacent README.md.
% Compile with XeLaTeX. Mathematical content is maintained in Markdown.
\documentclass[11pt,a4paper]{article}
\usepackage[fontsize=10.5pt]{scrextend}
\usepackage[margin=22mm,headheight=15pt,headsep=9mm,footskip=11mm]{geometry}
\usepackage{fontspec}
\setmainfont{texgyrepagella}[Extension=.otf,UprightFont=*-regular,BoldFont=*-bold,ItalicFont=*-italic,BoldItalicFont=*-bolditalic]
\setsansfont{texgyreheros}[Extension=.otf,UprightFont=*-regular,BoldFont=*-bold,ItalicFont=*-italic,BoldItalicFont=*-bolditalic]
\setmonofont{texgyrecursor}[Extension=.otf,UprightFont=*-regular,BoldFont=*-bold,ItalicFont=*-italic,BoldItalicFont=*-bolditalic,Scale=MatchLowercase]
\usepackage{amsmath,amssymb}
\usepackage{unicode-math}
\setmathfont{texgyrepagella-math.otf}
\usepackage{xcolor,microtype,enumitem,needspace,fancyhdr}
\usepackage{bookmark}
\definecolor{ink}{HTML}{17344B}
\definecolor{accent}{HTML}{197D80}
\definecolor{muted}{HTML}{596773}
\hypersetup{colorlinks=true,linkcolor=accent,urlcolor=accent,pdftitle={MI-26 - Concave
unitary-orbit subadditivity without
monotonicity},pdfauthor={Open Problems in Numerical Linear Algebra}}
\urlstyle{same}
\setlength{\parindent}{0pt}
\setlength{\parskip}{5pt plus 1pt minus 1pt}
\linespread{1.04}
\setlength{\emergencystretch}{3em}
\widowpenalty=10000
\clubpenalty=10000
\displaywidowpenalty=10000
\allowdisplaybreaks[1]
\setlist{nosep,leftmargin=1.5em}
\providecommand{\tightlist}{\setlength{\itemsep}{0pt}\setlength{\parskip}{0pt}}
\providecommand{\pandocbounded}[1]{#1}
\pagestyle{fancy}
\fancyhf{}
\fancyhead[L]{\sffamily\footnotesize\color{muted}OPEN PROBLEMS IN NUMERICAL LINEAR ALGEBRA}
\fancyhead[R]{\sffamily\bfseries\footnotesize\color{accent}MI-26}
\fancyfoot[L]{\parbox[b]{0.9\textwidth}{\sffamily\fontsize{8}{10}\selectfont\color{muted}Editorial ratings; a bounded search cannot certify that a problem remains unsolved.\\Literature check: 2026-09-12}}
\fancyfoot[R]{\sffamily\footnotesize\color{muted}\thepage}
\renewcommand{\headrulewidth}{0.3pt}
\renewcommand{\headrule}{\hbox to\headwidth{\color{accent}\leaders\hrule height \headrulewidth\hfill}}
\makeatletter
\renewcommand{\section}{\Needspace{5\baselineskip}\@startsection{section}{1}{0pt}{12pt plus 2pt minus 2pt}{4pt}{\sffamily\bfseries\large\color{ink}}}
\renewcommand{\subsection}{\Needspace{4\baselineskip}\@startsection{subsection}{2}{0pt}{9pt plus 1pt minus 1pt}{3pt}{\sffamily\bfseries\normalsize\color{ink}}}
\makeatother
\setcounter{secnumdepth}{0}
\begin{document}
{\sffamily\small\color{muted}Matrix inequalities and norms\par}
\vspace{3pt}
{\raggedright\hyphenpenalty=10000\exhyphenpenalty=10000\sffamily\bfseries\fontsize{21}{24}\selectfont\color{ink}Concave
unitary-orbit subadditivity without monotonicity\par}
\vspace{7pt}
{\sffamily\small\textbf{Difficulty:} challenging\quad\textbf{Importance:} interesting
to the community\par}
{\sffamily\small\bfseries\color{accent}Status: Lean verified\par}
\vspace{4pt}
{\color{accent}\hrule height 0.8pt}
\vspace{6pt}
\textbf{Rating rationale:} Removing monotonicity from a unitary-orbit
comparison is challenging; spectral bounds for broad classes of matrix
functions have community importance.

\subsection{Lean proof and verification evidence -
2026-09-12}\label{lean-proof-and-verification-evidence---2026-09-12}

\textbf{The complete original MI-26 assertion is false, with a
Lean-verified counterexample.} The real-valued concave function
\(f(x)=x-x^2\) and the rational projections

\[P=\begin{pmatrix}1&0\\0&0\end{pmatrix},\qquad
Q=\frac1{25}\begin{pmatrix}9&12\\12&16\end{pmatrix}\]

have \(f(P)=f(Q)=0\), while the genuine spectral matrix function
satisfies \(w^*f(P+Q)w=6/5>0\) for \(w=(1,-2)^{\mathsf T}\). Every pair
of complex unitary conjugates on the right is therefore zero. The
\href{https://github.com/sgstepaniants/OpenProblemsInNLA/tree/81176af27e570b59ba1e1a0745e28944e7d57c03/matrix-inequalities-and-norms/MI-26/lean}{proof
at revision 81176af} preserves the original complex PSD inputs and
complete real-valued function class.

\textbf{Mathematical counterexample and informal proof:} Matthew J.
Colbrook, Department of Applied Mathematics and Theoretical Physics,
University of Cambridge. \textbf{Lean formalization:} George
Stepaniants, Department of Computing and Mathematical Sciences,
California Institute of Technology, Pasadena, California, USA, with
AI-agent assistance.

The seven checked declarations in
\href{https://github.com/sgstepaniants/OpenProblemsInNLA/blob/81176af27e570b59ba1e1a0745e28944e7d57c03/matrix-inequalities-and-norms/MI-26/lean/Solution.lean}{Solution.lean}
are:

\begin{itemize}
\tightlist
\item
  \textit{NLA.MI26.admissibleFunction\_iff}: the exact original scalar
  concavity and value-at-zero conditions.
\item
  \textit{NLA.MI26.functionalCalculus\_eq\_spectral}: genuine CFC equals
  the finite spectral formula without a continuity premise.
\item
  \textit{NLA.MI26.functionalCalculus\_congr\_nonneg}: half-line
  extensions do not change PSD matrix functions.
\item
  \textit{NLA.MI26.quadratic\_cfc}: the witness function gives the
  actual matrix polynomial \(A-A^2\).
\item
  \textit{NLA.MI26.witness\_data}: rational projections, exact CFC
  values and the positive quadratic form.
\item
  \textit{NLA.MI26.counterexample}: the fixed pair defeats every pair of
  complex unitaries.
\item
  \textit{NLA.MI26.not\_subadditivityConjecture}: negation of the
  complete original assertion.
\end{itemize}

Two independent agents reviewed the
\href{https://github.com/ajt60gaibb/OpenProblemsInNLA/blob/main/matrix-inequalities-and-norms/MI-26/lean/reviews}{frozen
statements and completed proof}.
\href{https://github.com/sgstepaniants/OpenProblemsInNLA/actions/runs/34713045511}{Linux
run 34713045511} matched all seven exports with the sandboxed Comparator
and replayed the solution through Lean's default kernel. The
\href{https://github.com/ajt60gaibb/OpenProblemsInNLA/blob/main/matrix-inequalities-and-norms/MI-26/lean/verification/linux-2026-09-12}{original
artifacts and independent operational audit} retain all 118 source
hashes, isolation checks and rejection controls. All 15
\href{https://github.com/ajt60gaibb/OpenProblemsInNLA/blob/main/matrix-inequalities-and-norms/MI-26/lean/verification/proof-axioms.json}{transitive
axiom checks} use only \textit{propext}, \textit{Classical.choice} and
\textit{Quot.sound}. The explicit kernel LeanCert certificate \(0<6/5\)
remains in the final contradiction. These are independent agent reviews,
not external human peer review.

The proof pins \textbf{Lean 4.33.1},
\href{https://github.com/alerad/leancert/tree/621a43d7cf21f87872392a01e874f2f1dbddc926}{LeanCert
621a43d} and
\href{https://github.com/leanprover-community/mathlib4/tree/0df444a360eaa60ab8c11dca51a86af692955474}{Mathlib
0df444a}. See the
\href{https://github.com/ajt60gaibb/OpenProblemsInNLA/blob/main/matrix-inequalities-and-norms/MI-26/lean/formalization.yaml}{manifest},
\href{https://github.com/ajt60gaibb/OpenProblemsInNLA/blob/main/matrix-inequalities-and-norms/MI-26/lean/lake-manifest.json}{dependency
pins} and
\href{https://github.com/ajt60gaibb/OpenProblemsInNLA/blob/main/matrix-inequalities-and-norms/MI-26/lean/NUMERICAL_TARGETS.md}{numerical
targets}. On a documented
\href{https://github.com/ajt60gaibb/OpenProblemsInNLA/blob/main/tools/lean/HARNESS.md}{non-root
Linux host}, reproduce from the verified revision with:

\begin{verbatim}
tools/lean/bootstrap.sh /absolute/path/to/nla-lean-tools
tools/lean/selftest.sh /absolute/path/to/nla-lean-tools
tools/lean/verify.sh \
  matrix-inequalities-and-norms/MI-26/lean \
  /absolute/path/to/nla-lean-tools
\end{verbatim}

The optional positive-definite variant in the informal source is outside
these seven exports. The narrower globally nonnegative-valued function
class is not refuted. The complete original PSD target and its
historical informal proof remain below.

\subsection{Resolution --- 2026-09-11}\label{resolution-2026-09-11}

\textbf{Negative result by Matthew J. Colbrook}, Department of Applied
Mathematics and Theoretical Physics, University of Cambridge.
\textbf{Independent proof review: PASS.}

The real-valued concave function \(f(x)=x-x^2\) and two rational
projections refute the two-unitary inequality; an explicit positive
definite variant also works. The allowed condition is \(f(0)\ge0\),
without global nonnegativity or monotonicity. This does not refute the
narrower nonnegative-valued function class.

The exact target is resolved. The original statement and source evidence
are retained below; its former difficulty rating is historical.

\textbf{Primary manuscript:}
\href{https://github.com/ajt60gaibb/OpenProblemsInNLA/blob/main/matrix-inequalities-and-norms/MI-26/solution.pdf}{complete
proof PDF},
\href{https://github.com/ajt60gaibb/OpenProblemsInNLA/blob/main/matrix-inequalities-and-norms/MI-26/solution.tex}{standalone
TeX}, Theorem 1.1 and its proof;
\href{https://github.com/ajt60gaibb/OpenProblemsInNLA/blob/main/matrix-inequalities-and-norms/MI-26/solution.md}{authorship
and scope}. The
\href{https://github.com/ajt60gaibb/OpenProblemsInNLA/blob/main/references/colbrook-matrix-2026-09-11/verification/reviews/MI-26-review.md}{independent
review} checks the full original argument and records its hash. The
draft was AI-assisted; the 2026-09-11 review was independent agent
verification, without external human peer review or formal
certification. The later Lean verification above covers the complete
original PSD target.
\href{https://github.com/ajt60gaibb/OpenProblemsInNLA/blob/main/references/colbrook-matrix-2026-09-11/README.md}{Submission
record}.

\subsection{Problem statement}\label{problem-statement}

For every integer \(n\ge1\), positive semidefinite matrices
\(A,B\in\mathbb C^{n\times n}\), and real-valued concave function
\(f:[0,\infty)\to\mathbb R\) with \(f(0)\ge0\), must there exist unitary
matrices \(U,V\in\mathbb C^{n\times n}\) such that

\[f(A+B)\preceq Uf(A)U^*+Vf(B)V^*?\]

Here \(X\preceq Y\) means \(Y-X\) is positive semidefinite,
\(U^*U=UU^*=I\), and \(f(A)\) is defined by applying \(f\) to the
eigenvalues in a spectral decomposition of \(A\). Concavity means

\[f(\theta x+(1-\theta)y)\ge\theta f(x)+(1-\theta)f(y)
\quad (x,y\ge0,\ 0\le\theta\le1).\]

The function is not assumed nonnegative on its entire domain; adding
that assumption would change the problem. The unitaries may depend on
\(A,B,f\).

The requested inequality is an order comparison between sums of matrix
functions up to changes of orthonormal basis. Such comparisons yield
spectral and norm bounds for nonlinear transformations of positive
semidefinite matrices.

\subsection{References}\label{references}

\begin{enumerate}
\def\labelenumi{\arabic{enumi}.}
\tightlist
\item
  K. M. R. Audenaert and F. Kittaneh, \emph{Problems and Conjectures in
  Matrix and Operator Inequalities}, arXiv:1201.5232v3 (2012), §2,
  Problem 5 and equation (11).
  \href{https://arxiv.org/html/1201.5232}{Full text}.
\item
  J.-C. Bourin and E.-Y. Lee, \emph{Unitary orbits of Hermitian
  operators with convex or concave functions}, Bull. Lond. Math. Soc. 44
  (2012), 1085--1102, Theorem 3.1 and Remark 3.13.
  \href{https://arxiv.org/html/1109.2384}{Preprint}.
\end{enumerate}

Status check (2026-09-10): the two sources explicitly leave removal of
monotonicity open. The published 2017 Audenaert--Kittaneh update also
retains the question. Searches for ``concave unitary subadditivity
monotonicity'', ``Bourin Lee nonmonotone concave'', and 2025/2026,
together with the new arXiv:2609.05854 paper on Horn and orbit
inequalities, found no stated resolution of this full function class.
Recent results for nonnegative concave functions impose a stronger
assumption. This is a bounded status check.
\end{document}
